\documentclass[10pt]{amsart}
\usepackage[T1]{fontenc}
\usepackage{lmodern}
\usepackage[protrusion=true,expansion=false]{microtype}
\usepackage{mathtools,amssymb,mathrsfs,enumitem}
\usepackage[hidelinks]{hyperref}
\usepackage[nameinlink,capitalize,noabbrev]{cleveref}
\numberwithin{equation}{section}
\newtheorem{theorem}{Theorem}[section]
\newtheorem{proposition}[theorem]{Proposition}
\newtheorem{lemma}[theorem]{Lemma}
\newtheorem{corollary}[theorem]{Corollary}
\theoremstyle{remark}
\newtheorem{remark}[theorem]{Remark}
\newcommand{\CC}{\mathbf C}
\newcommand{\QQ}{\mathbf Q}
\newcommand{\PP}{\mathbf P}
\newcommand{\IC}{\operatorname{IC}}
\newcommand{\Alb}{\operatorname{Alb}}
\newcommand{\Stab}{\operatorname{Stab}}
\newcommand{\Sing}{\operatorname{Sing}}
\newcommand{\Sym}{\operatorname{Sym}}
\newcommand{\reg}{\mathrm{reg}}
\newcommand{\cO}{\mathcal O}
\newcommand{\cT}{\mathcal T}
\title[A realization of $F_4$ for Litt Problem No. 13]
{A Realization of $F_4$ for Litt Problem No. 13}
\author{OpenAI}
\date{September 30, 2026}
\subjclass[2020]{14K12, 14J30, 18M25, 20G41}
\keywords{Fano surface, Eckardt point, intersection complex, convolution,
Tannaka group, Albert algebra}
\hypersetup{pdftitle={A Realization of F4 for Litt Problem No. 13},pdfauthor={OpenAI}}
\begin{document}
\raggedbottom
\begin{abstract}
Let $F$ be the Fano surface of a smooth cubic threefold with an Eckardt
point, and let $E\subset F$ parametrize the lines through that point.
The image $X$ of $F$ in $\Alb(F)/E$ is an integral surface whose full
convolution Tannaka group is $F_4$ on its minimal $26$-dimensional module.
The quotient contracts precisely $E$ and splits the Albert module as
$27=1\oplus26$.  The fixed vector has nonzero Albert norm, as follows
from the Hessian ranks on the singular cubic and connected irreducibility.
The construction gives a seven-dimensional family of polarized ambient
fourfolds.  Its intersection complexes have characteristic cycle
$\Lambda_X+2\Lambda_0$ and generic twisted Hodge multiplicities $(6,14,6)$.
\end{abstract}
\maketitle

\section{Introduction}

For an integral subvariety $Z$ of a complex abelian variety, the
intersection complex $\IC_Z$ generates a Tannakian category under
convolution, after objects of Euler characteristic zero are discarded.
Litt Problem No.~13 asks whether the simple types $G_2$, $E_7$, $E_8$,
and $F_4$ occur as the resulting groups \cite{Litt13}.  The construction
below realizes $F_4$ by singular surfaces.

An Eckardt point of a smooth cubic threefold is a point whose tangent
hyperplane section is a cone over a smooth plane cubic.  Let $Y\subset\PP^4$
be such a threefold, with Eckardt point $p$.  Its Fano surface $F=F(Y)$
contains the elliptic curve $E$ of lines through $p$.  Choose an Albanese
embedding $F\hookrightarrow A=\Alb(F)$ based at a point of $E$, so that
$E\subset A$ is an abelian subvariety, and put
\[
 q:A\longrightarrow B=A/E,\qquad X=q(F)_{\mathrm{red}},\qquad f=q|_F.
\]

\begin{theorem}\label{thm:main}
For every smooth cubic threefold with a marked Eckardt point, $B$ is an
abelian fourfold and $X\subset B$ is an integral symmetric surface with
trivial translation stabilizer.  The map $f$ is birational, its only
positive-dimensional fibre is $f^{-1}(0)=E$, and
\[
 Rf_*\CC_F[2]\simeq\IC_X\oplus\delta_0,\qquad \chi(\IC_X)=26.
\]
The full convolution Tannaka pair is
\[
 \bigl(G(\IC_X),\omega(\IC_X)\bigr)\simeq(F_4,V_{26}),
\]
where $V_{26}$ is the minimal representation of $F_4$.
\end{theorem}

The proof uses the $E_6$ Tannaka group of the Fano surface.  Quotient
pushforward gives $H=G(\IC_X)\subset E_6$ and a restriction of the Albert
module $J|_H=1\oplus V_{26}$.  Triviality of the translation stabilizer
makes $V_{26}$ irreducible on $H^0$.  The Hessian of the Albert norm then
forces the fixed line to contain a norm-one vector.  Its full stabilizer
is $F_4$, and the classification of irreducible subgroups of the minimal
module gives $H=F_4$.

\section{Convolution and the normalized Fano surface}

All varieties are over $\CC$.  Intersection complexes have the perverse
normalization: if $j:Z_{\reg}\hookrightarrow Z$ and $d=\dim Z$, then
$\IC_Z=j_{!*}\CC_{Z_{\reg}}[d]$, extended by zero to the ambient abelian
variety.  Write $\delta_0$ for the skyscraper sheaf at the origin.  For
addition $m$ the convolution product is $K*L=Rm_*(K\boxtimes L)$.

The quotient by negligible objects is Tannakian on every rigid
subcategory generated by a semisimple perverse sheaf of geometric origin
\cite{KW2015,JKLM2025}.  Its group is reductive, and
$\dim\omega(K)=\chi(K)$.  A quotient homomorphism $q:A\twoheadrightarrow B$
induces a tensor functor and a closed immersion
\begin{equation}\label{eq:quotient-tensor}
 G(Rq_*K)\hookrightarrow G(K)
\end{equation}
\cite[Lemma~13.7]{KW2015}.  Groups refer throughout to the full group;
$E_6$ denotes its simply connected form.

The Albanese variety of $F$ has dimension five, and its Albanese maps
are embeddings \cite{CG1972}.  For $s\in F$ let $L_s\subset Y$ be its line.
Roulleau constructs an involution $\sigma=\sigma_E$ fixing $E$ pointwise
and a map $\gamma_E:F\to E$ such that, for general $s\notin E$,
\begin{equation}\label{eq:plane-section}
 Y\cap\langle p,L_s\rangle=L_s+L_{\gamma_E(s)}+L_{\sigma(s)}.
\end{equation}
The induced Albanese automorphism has eigenvalue $+1$ on $T_0E$ and
$-1$ on $T_0A/T_0E$ \cite[Section~3.1 and Theorem~13]{Roulleau2009}.
Thus $q(\sigma(s))=-q(s)$ and $X=-X$.

\begin{lemma}\label{lem:normalization}
The Albanese embedding can be chosen with $E\subset A$ an abelian
subvariety, the sum of every coplanar triple equal to zero, and
$G(\IC_F)=E_6$ on the Albert module $J$.
\end{lemma}
\begin{proof}
For an embedding based at $e\in E$, addition is constant on the incidence
variety $\cT\subset F^3$ of ordered triples forming a plane section
\cite{CG1972,Kraemer2022}; call the value $c\in A$.  By
\eqref{eq:plane-section}, its image in $B$ is $x+0-x=0$, so $c\in E$.
Choose $u\in E$ with $3u=c$ and translate the embedding by $-u$.
This preserves $E$ and the quotient and makes the triple sum zero.

The incidence correspondence now gives a nonzero invariant symmetric
cubic on $J$ \cite[Example~2.2(c)]{Kraemer2022}.  The derived connected
Tannaka group and its representation are $(E_6,J)$
\cite[Theorem~2]{Kraemer2016}.  On this module the invariant cubic is a
nonzero multiple of the Albert norm, whose full linear stabilizer is
$E_6$ \cite[Theorem~2.1(c)]{Kraemer2022}.  Faithfulness therefore gives
$E_6\subset G(\IC_F)\subset E_6$.
\end{proof}

\section{The quotient geometry}

The tangent bundle theorem identifies
\begin{equation}\label{eq:tangent}
 \PP(T_sF)=L_s\subset\PP(T_0A)=\PP^4,
 \qquad \PP(T_0E)=\{p\}
\end{equation}
\cite{CG1972}\cite[Proposition~4]{Roulleau2009}.

\begin{proposition}\label{prop:birational}
The map $f:F\to X$ is birational and immersive on $F\setminus E$.
Its only contracted curve is $E$, and $f^{-1}(0)=E$.
\end{proposition}
\begin{proof}
By \eqref{eq:tangent}, $\ker(df_s)\ne0$ exactly when $p\in L_s$, that
is, when $s\in E$.  Any contracted curve consequently lies in $E$.
Conversely, $E$ is contracted, and $q^{-1}(0)=E$ in $A$ gives the stated
fibre.  Thus $f$ is generically finite onto an integral surface.

Choose a dense open $U\subset X_{\reg}$ where $f$ is finite and unramified,
avoiding $B[2]$ and the nongeneric locus of \eqref{eq:plane-section}.
For $s,t\in f^{-1}(x)$ with $x\in U$, projection from $p$ sends both
$L_s$ and $L_t$ to $\PP(T_xX)$.  Hence both lie in
$\Pi=\langle p,L_s\rangle$.  The plane section
\eqref{eq:plane-section} has exactly three lines.  Its line through $p$
maps to zero, so $t=s$ or $t=\sigma(s)$.  The latter implies
$x=-x$, contrary to $x\notin B[2]$.  The generic degree is therefore one.
\end{proof}

\begin{proposition}\label{prop:semismall}
One has $Rf_*\CC_F[2]\simeq\IC_X\oplus\delta_0$ and $\chi(\IC_X)=26$.
\end{proposition}
\begin{proof}
The proper birational surface map $f$ is semismall.  Its only fibre of
positive dimension is the irreducible curve $E$ over zero.
Roulleau computes $E^2=-3$ \cite[Proposition~10]{Roulleau2009}.
The local intersection form is the nondegenerate matrix $(-3)$, so the
semismall decomposition has exactly one skyscraper summand
\cite[Theorem~3.4.1]{deCM2002}.  Since $e(F)=27$ \cite{CG1972},
the Euler characteristic of the remaining summand is $27-1=26$.
\end{proof}

\begin{proposition}\label{prop:stabilizer}
$\Stab_B(X)=\{0\}$, and the tautological module of $G(\IC_X)$ is
irreducible on the identity component.
\end{proposition}
\begin{proof}
Factor $f$ through the normalization $\nu:X^\nu\to X$, with
$h:F\to X^\nu$.  The only exceptional curve of $h$ is $E$.  A proper
birational morphism of normal surfaces is an isomorphism outside its
exceptional locus, so $X^\nu\setminus\{z\}\simeq F\setminus E$, where
$z=h(E)$.  The point $z$ is singular: a birational morphism to a smooth
surface factors locally into point blowups, whose exceptional curves are
rational, whereas $E$ is elliptic.  Thus $\Sing(X^\nu)=\{z\}$.

A translation preserving $X$ lifts to its normalization and fixes $z$.
Since $\nu(z)=0$, the translation is zero.  This is precisely
nondivisibility of the simple, non-negligible sheaf $\IC_X$.
Connected irreducibility follows from
\cite[Proposition~3.3(2)]{JKLM2025}.
\end{proof}

\section{The Albert norm and the full group}

Put $H=G(\IC_X)$ and $V=\omega(\IC_X)$.  By
\cref{lem:normalization,prop:semismall} and \eqref{eq:quotient-tensor},
\begin{equation}\label{eq:restriction}
 H\hookrightarrow E_6,\qquad J|_H\simeq1\oplus V,\qquad\dim V=26.
\end{equation}
Indeed the tensor unit already belongs to $\langle\IC_X\rangle$, so
$\langle\IC_X\oplus\delta_0\rangle=\langle\IC_X\rangle$.
Let $v_0\ne0$ span the trivial summand, and let $N$ be the Albert cubic.
Its complete polarization defines $\beta_v:J\to J^\vee$ by
$\beta_v(u)=N(v,u,-)$.

\begin{lemma}\label{lem:hessian}
For $v\ne0$ with $N(v)=0$, the kernel of $\beta_v$ has dimension $17$
in Albert rank one and $9$ in Albert rank two.
\end{lemma}
\begin{proof}
In octonion coordinates the cubic is
\[
 N=abc-a\,n(\xi)-b\,n(\eta)-c\,n(\zeta)
       +2\operatorname{Re}((\xi\eta)\zeta),
\]
where each octonion norm is nondegenerate of rank eight.
The rank-orbit description of the Albert algebra reduces the two cases
to $e_1=\operatorname{diag}(1,0,0)$ and $e_1+e_2$
\cite{SpringerVeldkamp2000}.  The quadratic terms at these points are,
respectively,
\[
 bc-n(\xi),\qquad \gamma(\alpha+\beta)-n(\xi)-n(\eta).
\]
Their ranks are $10$ and $18$, giving kernel dimensions $27-10$ and $27-18$.
\end{proof}

\begin{proposition}\label{prop:fixed-norm}
The fixed vector satisfies $N(v_0)\ne0$.
\end{proposition}
\begin{proof}
The map $\beta_{v_0}$ is $H^0$-equivariant.  By
\cref{prop:stabilizer}, $J|_{H^0}=1\oplus V$ is a multiplicity-free sum
of two irreducibles.  Reductivity therefore gives
\[
 \dim\ker\beta_{v_0}\in\{0,1,26,27\}.
\]
If $N(v_0)=0$, the vector has Albert rank one or two, contradicting
\cref{lem:hessian}.
\end{proof}

Scale $v_0$ to norm one.  Its full stabilizer in $E_6$ is the connected
group $F_4$ \cite{SpringerVeldkamp2000}.  Hence $H\subset F_4$, and $V$
is its minimal module.  The group $H^0$ is positive-dimensional, since it
acts irreducibly on a space of dimension $26$.  No proper
positive-dimensional closed subgroup of $F_4$ in characteristic zero acts
irreducibly on this module: this is the $F_4$ row of
\cite[Theorem~1 and Table~1.2]{LiebeckSeitz2004}.  Thus
\[
 F_4=H^0\subset H\subset F_4.
\]
This proves \cref{thm:main}, including connectedness of the full group.

\section{The family and its numerical invariants}

\begin{corollary}\label{cor:family}
The construction yields a seven-dimensional family of polarized ambient
fourfolds of type $(1,2,2,2)$.  Every member carries the surface of
\cref{thm:main}.
\end{corollary}
\begin{proof}
For a smooth cubic surface $S=\{g=0\}\subset\PP^3$ and a transverse
hyperplane $\{\ell=0\}$, the cubic
$Y=\{g+x_0^2\ell=0\}\subset\PP^4$ is smooth and has Eckardt point
$[1:0:0:0:0]$.  Let $P=\operatorname{Im}(1-\sigma)\subset J(Y)=A$.
The invariant part is $E$, and the principal polarization identifies
$A/E$ with $P^\vee$.  The polarization on $P$ has type $(1,1,1,2)$;
its dual has type $(1,2,2,2)$
\cite[Lemma~1.12 and Remark~1.14]{CMZ2021}.
The moduli space of these cubic surface pairs has dimension seven, and
its period map to polarized $P$ is injective with injective differential
\cite[Theorem~0.1]{CMZ2021}.  Dualization gives the asserted dimension
for the family of $B$.
\end{proof}

\begin{corollary}[A concrete example]\label{cor:example}
The cubic $x_0^2x_1+x_1^3+x_2^3+x_3^3+x_4^3=0$, marked at
$p=[1:0:0:0:0]$, gives a realization of $(F_4,V_{26})$.
\end{corollary}
\begin{proof}
The partial derivatives are $2x_0x_1$, $x_0^2+3x_1^2$, and $3x_i^2$
for $i=2,3,4$; they have no common projective zero.  The tangent section
at $p$ is the cone $x_2^3+x_3^3+x_4^3=0$ in $x_1=0$.
Apply \cref{thm:main}.
\end{proof}

For a subvariety $Z\subset B$, write $\Lambda_Z$ for the closure of the
conormal bundle to $Z_{\reg}$ in $T^*B$.

\begin{proposition}\label{prop:cycle}
The characteristic cycle and conormal Gauss degree are
\[
 \operatorname{CC}(\IC_X)=\Lambda_X+2\Lambda_0,
 \qquad \deg(\Lambda_X\to T_0^*B)=24.
\]
\end{proposition}
\begin{proof}
Over $B\setminus\{0\}$, $f$ is finite and immersive.  Locally its image
is a union of smooth branches, and its direct image is the sum of their
constant intersection complexes.  By \cref{prop:semismall}, no other
conormal component occurs away from zero.  Thus
$\operatorname{CC}(\IC_X)=\Lambda_X+m\Lambda_0$.

A general covector in $T_0^*B\subset T_0^*A$ defines a general hyperplane
$H\subset\PP^4$ through $p$.  By \eqref{eq:tangent}, a conormal point
away from zero is a line $L_s\subset Y\cap H$ not passing through $p$.
The surface $Y\cap H$ is a smooth cubic surface with Eckardt point $p$.
Of its $27$ lines, exactly three pass through $p$.  The remaining $24$
give the generic Gauss fibre; exceptional strata are avoided by a general
covector.  Hence the degree is $24$.  The index formula on abelian
varieties \cite[Theorem~1.3]{FK2000} now gives $26=24+m$.
\end{proof}

\begin{proposition}\label{prop:hodge}
For a general unitary character $\alpha$ of $B$, the twisted
hypercohomology $H^0(B,\IC_X\otimes\mathcal L_\alpha)$ is pure of weight
two with Hodge multiplicities $(6,14,6)$.
\end{proposition}
\begin{proof}
Generic vanishing and \cref{prop:semismall} give
$H^i(F,f^*\mathcal L_\alpha)=0$ for $i\ne2$ and dimension $27$ in
degree two \cite[Theorem~1.1]{KW2015}.  The unitary Hodge decomposition
therefore makes the higher cohomology of
$K_F\otimes f^*\alpha$ vanish.  Riemann--Roch gives
$h^0(K_F\otimes f^*\alpha)=\chi(K_F)=6$, since $p_g(F)=10$ and
$q(F)=5$ \cite{CG1972}.  The same calculation for $\alpha^{-1}$ and
duality give Hodge multiplicities $(6,15,6)$ for the twisted degree-two
cohomology of $F$.

The rational Hodge module form of the semismall decomposition is
$Rf_*\QQ_F^H[2]=\IC_X^H\oplus\QQ_0^H(-1)$
\cite{deCM2002,Saito1988}.  The last term has type $(1,1)$, yielding
$(6,14,6)$ after its removal.
\end{proof}

\begin{remark}
The $24$ short roots and the zero weight of multiplicity two in $V_{26}$
agree with the two terms of the characteristic cycle.  The numerical
calculation is independent of the subgroup argument.
\end{remark}

\begin{thebibliography}{99}
\bibitem{CG1972} C.~H. Clemens and P.~A. Griffiths,
\emph{The intermediate Jacobian of the cubic threefold},
Ann. of Math. (2) \textbf{95} (1972), 281--356.
\bibitem{CMZ2021} S. Casalaina-Martin and Z. Zhang,
\emph{The moduli space of cubic surface pairs via the intermediate
Jacobians of Eckardt cubic threefolds}, J. Lond. Math. Soc. (2)
\textbf{104} (2021), no.~1, 1--34.
\bibitem{deCM2002} M.~A.~A. de Cataldo and L. Migliorini,
\emph{The hard Lefschetz theorem and the topology of semismall maps},
Ann. Sci. \'Ecole Norm. Sup. (4) \textbf{35} (2002), no.~5, 759--772.
\href{https://doi.org/10.1016/S0012-9593(02)01108-4}{doi:10.1016/S0012-9593(02)01108-4}.
\bibitem{FK2000} J. Franecki and M. Kapranov,
\emph{The Gauss map and a noncompact Riemann--Roch formula for
constructible sheaves on semiabelian varieties}, Duke Math. J.
\textbf{104} (2000), no.~1, 171--180.
\bibitem{JKLM2025} A. Javanpeykar, T. Kr\"amer, C. Lehn and M. Maculan,
\emph{The monodromy of families of subvarieties on abelian varieties},
Duke Math. J. \textbf{174} (2025), no.~6, 1045--1149.
\href{https://doi.org/10.1215/00127094-2024-0053}{doi:10.1215/00127094-2024-0053}.
\bibitem{Kraemer2016} T. Kr\"amer,
\emph{Cubic threefolds, Fano surfaces and the monodromy of the Gauss map},
Manuscripta Math. \textbf{149} (2016), 303--314.
\href{https://doi.org/10.1007/s00229-015-0785-z}{doi:10.1007/s00229-015-0785-z}.
\bibitem{Kraemer2022} T. Kr\"amer,
\emph{Characteristic cycles and the microlocal geometry of the Gauss map, I},
Ann. Sci. \'Ecole Norm. Sup. (4) \textbf{55} (2022), no.~6, 1475--1527.
\href{https://doi.org/10.24033/asens.2522}{doi:10.24033/asens.2522}.
\bibitem{KW2015} T. Kr\"amer and R. Weissauer,
\emph{Vanishing theorems for constructible sheaves on abelian varieties},
J. Algebraic Geom. \textbf{24} (2015), no.~3, 531--568.
\href{https://doi.org/10.1090/jag/645}{doi:10.1090/jag/645}.
\bibitem{LiebeckSeitz2004} M.~W. Liebeck and G.~M. Seitz,
\emph{Subgroups of exceptional algebraic groups which are irreducible on
an adjoint or minimal module}, J. Group Theory \textbf{7} (2004), no.~3,
347--372.
\href{https://doi.org/10.1515/jgth.2004.012}{doi:10.1515/jgth.2004.012}.
\bibitem{Litt13} D. Litt, \emph{Problem No.~13}, Problems I Like,
\url{https://www.problemsilike.com/13}, accessed September 30, 2026.
\bibitem{Roulleau2009} X. Roulleau,
\emph{Elliptic curve configurations on Fano surfaces},
Manuscripta Math. \textbf{129} (2009), 381--399.
\href{https://doi.org/10.1007/s00229-009-0264-5}{doi:10.1007/s00229-009-0264-5}.
\bibitem{Saito1988} M. Saito, \emph{Modules de Hodge polarisables},
Publ. Res. Inst. Math. Sci. \textbf{24} (1988), no.~6, 849--995.
\bibitem{SpringerVeldkamp2000} T.~A. Springer and F.~D. Veldkamp,
\emph{Octonions, Jordan Algebras and Exceptional Groups},
Springer Monographs in Mathematics, Springer, Berlin, 2000.
\href{https://doi.org/10.1007/978-3-662-12622-6}{doi:10.1007/978-3-662-12622-6}.
\end{thebibliography}
\end{document}
